\section{A two-bit reverse dispatch removes state 223}
\label{sec:fold236}

Let $B$ be the 237-state table obtained by the certified closed-chunk
reassociation named \texttt{n234\_left}, and let $C$ be the delivered
236-state table. All state numbers in this section are inherited labels;
the start and halt labels are 95 and 254. The complete change from $B$ to
$C$ is
\[
\begin{array}{c|c|c}
\text{row}&B&C\\\hline
182/0&0\,R\,187&0\,L\,198\\
182/1&1\,R\,212&1\,L\,210\\
237/0&0\,R\,223&0\,R\,182.
\end{array}
\]
The two rows of state 223 are then deleted. Define $\phi(223)=182$ and
$\phi(q)=q$ for every other state. On configurations, $\phi$ changes only
the control state and preserves the tape and head position.

\paragraph{The guard is a predecessor theorem.}
Inspection of every row of the literal old table gives exactly one
incoming row for each relevant state:
\[
76/1\longmapsto 1\,R\,182,
\qquad
237/0\longmapsto 0\,R\,223.
\]
Since the initial state is 95, induction on elapsed old-machine time
proves the following statement on \emph{every binary initial tape}:
if $B$ is in state 182, the cell immediately left of its head contains 1;
if it is in state 223, that cell contains 0. No register values, source
program counter, run-length restriction, finite support assumption, or
empirical coverage assumption enters this induction.

\paragraph{Four exact blocks.}
Write $a$ for the bit left of the head and $b$ for the bit under the head.
The new state 182 first reads $b$, moves left without changing it, and
chooses state 198 when $b=0$ or state 210 when $b=1$. That helper reads
$a$, moves right without changing it, and chooses a second helper. The
second helper rereads $b$ and completes the old transition:
\[
\begin{array}{c|c|c|l}
\text{old state}&b&a&\text{new three-step state path}\\\hline
182&0&1&182\to198\to136\to187\\
182&1&1&182\to210\to177\to212\\
223&0&0&182\to198\to206\to23\\
223&1&0&182\to210\to114\to70.
\end{array}
\]
Each block reads only the cells at offsets $-1$ and $0$, never changes
any tape bit, and ends one cell right of the original head. Its final
state is exactly the image under $\phi$ of the old one-step final state.
Every other old row---470 rows in total---commutes with $\phi$ in one
step. This includes the redirected incoming row $237/0$.

The standalone checker \texttt{check\_local\_fold.py} reconstructs the
incoming-edge lists, compares all 470 ordinary rows, and computes these
four paths directly from the two literal tables on all 32 five-cell
windows satisfying the required tag. It rejects any read outside the
window and checks equality of the entire window, head position, and
mapped endpoint. The extra three window bits are arbitrary; the paths
never inspect them, so the same statement holds for an arbitrary infinite
tape outside the window.

\paragraph{Unbounded simulation and halting equivalence.}
Let $q_i$ denote the state of $B$ immediately before its $i$th transition,
and set
\[
 A(t)=\#\{i<t:q_i\in\{182,223\}\},\qquad
 \tau(t)=t+2A(t).
\]
Induction using the four blocks and the ordinary commuting rows gives
\[
 C^{\tau(t)}(c_0)=\phi\bigl(B^t(c_0)\bigr),
 \qquad t\leq\tau(t)\leq3t,
\]
for every initial configuration $c_0$ with state 95, while the old run is
defined. Every simulated block has positive length, and the four expanded
blocks contain no halt state. The map fixes both the start and halt
states. Consequently $B$ halts from a given initial tape if and only if
$C$ halts from that same tape. This is an ordinary unbounded mathematical
proof, formalizable by elementary induction; the finite checker verifies
its complete local premises rather than sampling a time horizon.

\paragraph{Why the source rotation matters.}
In the original composed237 table the row $177/1$ enters state 185. The
closed-chunk reassociation changes that row to $1\,R\,212$. This creates
the missing fourth route in the reverse dispatch above. The source
transformation itself keeps the number of states at 237, but exposes a
state sharing that the unchanged table does not have. The saving is
therefore a composition of a source-equivalent change and a local binary
encoding change; it is not merely a repeat of the three already composed
savings.

Relative to the original composed237 table, the final machine deletes
state 223, changes seven retained rows, and leaves the other 465 retained
rows unchanged. The source transformation's separate certificate and the
blank-start/runtime calibration complete the link to ZFC consistency.
The present local proof alone asserts halting equivalence of $B$ and $C$;
it makes no independent consistency assumption or proof-assistant claim.
